\section{Needs, features and working hypotheses}
\label{sec:specifications}

\subsection{Needs addressed by the application}

We designed NEO-Watch for two kinds of users. The first are \emph{astronomy enthusiasts}: they
want to know which objects come close to the Earth, to follow some of them over time, and to
be warned of the passages that interest them without checking the data every day. The second
are the \emph{administrators} of the application: they manage the accounts and make sure that
the data stays up to date.

\medskip

The NASA API alone does not meet these needs. It answers one request at a time, it has no
notion of user, of favourite or of alert, and its main endpoint, \texttt{/feed}, returns at most
seven days of data per call. NEO-Watch therefore stores the data once in its own database, and
builds the user features on top of it.

\subsection{Features delivered}

The subject required six features (F1 to F6). We also chose three optional features
(FO3, FO7, FO8) in the analysis report. All of them are implemented.
Table~\ref{tab:features} summarises what each feature does in the final application.

\begin{table}[H]
  \centering
  \small
  \begin{tabularx}{\textwidth}{@{}l>{\raggedright\arraybackslash}X l@{}}
    \toprule
    \textbf{Feature} & \textbf{What the user can do} & \textbf{Status} \\
    \midrule
    F1 Accounts & Create an account, log in (JWT tokens). Administrators list, promote and delete
      accounts, and see the connection history of a user. & Done \\
    F2 Dashboard & See the approaches of a chosen period (next 7 days by default), the number of
      NEOs and of hazardous NEOs, the closest approach, and their favourites. & Done \\
    F3 Search & Search by name, identifier, diameter, hazard flag, origin and approach date; sort
      the results; see the details and approaches of a NEO (km and astronomical units); export
      to CSV. & Done \\
    F4 Declare a NEO & Record an object they observed. It is stored with the origin
      \texttt{USER} and is never overwritten by the NASA imports. & Done, simplified \\
    F5 Favourites & Add and remove favourites, see their approaches, plot their distance to the
      Earth over time, export the history (CSV) and the plot (PNG). & Done \\
    F6 Alerts & Define alerts (minimum diameter, maximum distance, hazardous only), switch them
      on and off, and be notified once of each matching approach when logging in. & Done \\
    FO3 Statistics & Approaches per month, size distribution, share of hazardous NEOs, mean
      and median distance. & Done \\
    FO7 Automatic import & The NASA data is updated automatically every day; administrators
      can also start an update by hand. & Done \\
    FO8 Picture of the day & See and download NASA's Astronomy Picture of the Day, for any
      date. & Done \\
    \bottomrule
  \end{tabularx}
  \caption{Features of NEO-Watch and their status}
  \label{tab:features}
\end{table}

\subsection{Working hypotheses}

The application relies on a few hypotheses, decided in the analysis report and kept during
the development.

\begin{itemize}
  \item \textbf{Data source.} NEO-Watch only uses the \texttt{/feed} endpoint of NeoWs. It is
    organised by period, like our features (dashboard, search by date, alerts), and filling three
    years of data takes about 160 calls, against more than 3,000 with \texttt{/neo/browse}. As a
    consequence, the application only knows the NEOs that come close to the Earth during its
    import window: it monitors approaches, it is not a full catalogue.
  \item \textbf{Import window.} Two years in the past and one year in the future. The past is
    imported once. Every day, the application adds the missing days at the end of the window and
    imports the next 30 days again, to get the objects discovered since the last import.
  \item \textbf{Units.} Distances and diameters in kilometres, speeds in kilometres per hour.
    Distances in astronomical units are computed when they are displayed, not stored.
  \item \textbf{NASA data is read-only.} The application never changes a NASA object; users only
    add their own objects (F4), their favourites and their alerts.
  \item \textbf{Notifications at login.} Alerts are evaluated when the user logs in. Each
    matching approach is notified once.
  \item \textbf{Accounts.} The first administrator is created with the seed data; every other
    administrator is promoted by an existing one.
\end{itemize}

\subsection{Differences with the analysis report}

Some elements of the analysis report were changed or dropped during the development.
Table~\ref{tab:changes} lists them, with the reason for each change.

\begin{table}[H]
  \centering
  \small
  \begin{tabularx}{\textwidth}{@{}>{\raggedright\arraybackslash}X>{\raggedright\arraybackslash}X>{\raggedright\arraybackslash}X@{}}
    \toprule
    \textbf{Planned} & \textbf{Implemented} & \textbf{Reason} \\
    \midrule
    User NEOs and corrections of NEO fields submitted as \emph{proposals}, accepted or rejected
      by an administrator (\texttt{proposal} table). &
    Users create their NEOs directly, with the origin \texttt{USER}. The \texttt{proposal} table
      exists but is not used. &
    We gave priority to the six required features. The origin column still separates user
      objects from NASA objects. \\
    \addlinespace
    Call \texttt{/neo/\{id\}} when a NEO becomes a favourite, to get its full approach history. &
    The history of a favourite comes from the approaches already imported. &
    Fewer API calls and simpler code. The drawback: an object that passes rarely has only one or
      two points on its plot. \\
    \addlinespace
    Use the minimum estimated diameter for filtering, to be cautious in alerts (an alert ``larger
      than 100 m'' only for objects surely larger than 100 m). &
    The alert compares its threshold with the \emph{maximum} estimated diameter. &
    No object that may be large enough is missed, but an alert can fire for an object that is
      smaller than the threshold. This point should be harmonised. \\
    \addlinespace
    An alert on a specific NEO chosen by the user. &
    Not implemented. &
    Following a given NEO is already the role of the favourites (F5). \\
    \addlinespace
    A \texttt{search} table recording the searches. &
    Created but not used. &
    It would only have served usage statistics, which no feature needed. \\
    \addlinespace
    A session opened at login (\texttt{/api/v1/auth/login}). &
    Signed JWT tokens: an access token (30 minutes) and a refresh token (7 days). &
    A JWT is verified by its signature, without storing anything on the server; this solution
      was suggested to us during the project follow-up. \\
    \addlinespace
    Personal endpoints under \texttt{/users/me/} (the user is taken from the session). &
    Endpoints with the user identifier in the path (\texttt{/favourites/user/\{user\_id\}}), checked
      against the token: a user can only reach their own data (error 403 otherwise). &
    Simpler for the administrators, who can reach any user's data; the check on the token keeps
      the same protection. \\
    \bottomrule
  \end{tabularx}
  \caption{Differences between the analysis report and the final application}
  \label{tab:changes}
\end{table}
